% !TEX root = ../notes_template.tex
\chapter{Cervical \& Thoracic}\label{chp:cervical_thoracic }

\section*{Overview}

This week includes two discrete skill labs: one dedicated to cervical spine muscle testing and deep neck flexor evaluation, and the other to comprehensive upper and lower quarter neurological screens. These assessments provide foundational tools for screening neurologic integrity and motor control relevant to both orthopedic and neurorehabilitation contexts. No palpation or RUSI is included this week.

\section*{Objectives}
\begin{itemize}
  \item Apply standardized MMT protocols for cervical flexion, extension, rotation, and sidebending.
  \item Identify primary cervical muscles, their innervation, and spinal root levels.
  \item Execute upper and lower quarter screens using dermatomes, myotomes, and reflexes.
  \item Interpret screen findings to differentiate between spinal and peripheral nervous system involvement.
  \item Recognize clinical indications for a full neurological examination or referral.
\end{itemize}

\section*{Cervical Spine Muscle Testing}

\subsection*{Functional Context}

Cervical musculature plays a key role in posture, head positioning, and motor control. Differentiating between upper and lower cervical function allows precise testing of motor performance and neuromuscular coordination.

\subsection*{Cervical Movements: Muscles \& Innervation}

\begin{table}[H]
\centering
\captionsetup{justification=raggedright, singlelinecheck=false}
\begin{tabular}{|p{4cm}|p{5cm}|p{4cm}|p{2.5cm}|}
\hline
\textbf{Movement} & \textbf{Muscle} & \textbf{Peripheral Nerve} & \textbf{Nerve Root(s)} \\
\hline
\multirow{2}{*}{Capital Flexion}
  & Longus Capitis & C1–C3 Ventral Rami & C1–C3 \\
  & Rectus Capitis Anterior & C1 Ventral Ramus & C1 \\
\hline
\multirow{2}{*}{Capital Extension}
  & Rectus Capitis Posterior Major/Minor & Suboccipital & C1 \\
  & Obliquus Capitis Superior & Suboccipital & C1 \\
\hline
\multirow{2}{*}{Cervical Flexion}
  & Sternocleidomastoid & Spinal Accessory (CN XI), C2–C3 & CN XI, C2–C3 \\
  & Longus Colli & C2–C6 Ventral Rami & C2–C6 \\
\hline
\multirow{2}{*}{Cervical Extension}
  & Semispinalis Cervicis & Dorsal Rami of Cervical Spinal Nerves & C2–C8 \\
  & Cervical Erector Spinae & Dorsal Rami of Cervical Spinal Nerves & C3–C8 \\
\hline
\multirow{2}{*}{Lateral Flexion}
  & Upper Trapezius & Spinal Accessory (CN XI), C3–C4 & CN XI, C3–C4 \\
  & Scalenes (Ant/Mid/Post) & C3–C8 Ventral Rami & C3–C8 \\
\hline
\multirow{2}{*}{Rotation}
  & Sternocleidomastoid (contralateral) & Spinal Accessory (CN XI), C2–C3 & CN XI, C2–C3 \\
  & Splenius Capitis (ipsilateral) & Dorsal Rami of Cervical Spinal Nerves & C3–C5 \\
\hline
\end{tabular}
\caption{Cervical Spine Movements and Associated Neuromuscular Innervation}
\end{table}

\subsection*{MMT – Cervical Spine}

Manual muscle testing in the cervical spine requires a different approach than testing the limbs or even the lumbar trunk. While limb MMT emphasizes peak force and large joint excursions, and trunk MMT uses gross functional movements (e.g., clearing scapulae or sternum), cervical MMT emphasizes motor control, alignment, and segmental coordination.

These tests:

\begin{itemize}
  \item Highlight the ability to maintain cervical alignment and symmetry under low-force conditions.
  \item Require refined stabilization of the trunk, scapulae, and occiput to reduce substitution from superficial muscles like the sternocleidomastoids or upper trapezius.
  \item Depend on precise positioning, given short lever arms and high segmental sensitivity.
  \item Are interpreted in the context of postural patterns, pain inhibition, or neurological compromise rather than pure strength deficits.
\end{itemize}

Compared to lumbar MMT, cervical testing involves smaller ranges, increased sensitivity to substitution, and greater emphasis on segmental motor integrity and control.

\begin{labactivitybox}
\textbf{Activity: MMT – Capital Flexion and Extension}

\begin{itemize}
  \item \textbf{Capital Flexion:} Test \textit{longus capitis} and \textit{rectus capitis anterior}. Patient performs a subtle nod (“yes” motion) without lifting the head from the table.
  \item \textbf{Capital Extension:} Test \textit{rectus capitis posterior major/minor} and \textit{obliquus capitis superior}. Patient gently presses occiput into examiner’s hand without cervical spine involvement.
  \item Apply graded resistance and assess for compensatory strategies (e.g., SCM activation).
\end{itemize}
\end{labactivitybox}

\begin{labactivitybox}
\textbf{Activity: MMT – Lower Cervical Flexion and Extension}

\begin{itemize}
  \item \textbf{Cervical Flexion:} Patient flexes neck by tucking chin and lifting head off the table. Primary muscles: \textit{SCM}, \textit{longus colli}.
  \item \textbf{Cervical Extension:} Patient extends lower cervical spine off the table. Primary muscles: \textit{cervical erector spinae}, \textit{semispinalis cervicis}.
  \item Grade resistance and observe segmental motion and muscle balance.
\end{itemize}
\end{labactivitybox}

\begin{labactivitybox}
\textbf{Activity: MMT – Cervical Lateral Flexion and Rotation}

\begin{itemize}
  \item \textbf{Lateral Flexion:} Patient sidebends toward shoulder without rotation. Key muscles: \textit{SCM}, \textit{upper trapezius}, \textit{scalenes}.
  \item \textbf{Rotation:} Test both sides. SCM acts contralaterally; splenius capitis acts ipsilaterally.
  \item Assess motion symmetry, strength, and substitution strategies.
\end{itemize}
\end{labactivitybox}

\subsection*{Deep Neck Flexor Testing}

\begin{labactivitybox}
\textbf{Activity: DNF Endurance Test}

\begin{itemize}
  \item Position: Supine, chin gently tucked.
  \item Goal: Hold cranio-cervical flexion as long as possible without compensation.
  \item Observe for loss of form or chin jutting.
\end{itemize}
\end{labactivitybox}

\begin{labactivitybox}
\textbf{Activity: Cranio-Cervical Flexion Test (CCFT)}

\begin{itemize}
  \item Use pressure biofeedback inflated to 20 mmHg.
  \item Cue subtle chin tuck to gradually reach 22–30 mmHg in 2 mmHg increments.
  \item Hold each stage 10 seconds without substitution.
  \item Document motor control and precision.
\end{itemize}
\end{labactivitybox}




\section*{Upper and Lower Quarter Screens}

Quarter screens are rapid neurological screening tools designed to detect patterns of peripheral nervous system involvement. These screens use a combination of myotomal strength testing, dermatomal sensory checks, and deep tendon reflexes to assess spinal nerve root integrity. While not a substitute for detailed neurodiagnostic testing, quarter screens provide essential information for triage, differential diagnosis, and referral decisions—especially in cases of extremity pain with potential spinal origins. Their use supports efficient, systems-based examination in both orthopedic and neurologic contexts.

\subsection*{Upper Quarter Screen}

\begin{itemize}
  \item \textbf{Dermatomes:} C3–T2
  \item \textbf{Myotomes:} C5–T1
  \item \textbf{Reflexes:} C5 (biceps), C6 (brachioradialis), C7 (triceps)
  \item \textbf{MMT Examples:} Deltoid (C5), Triceps (C7), Thumb Ext (C8)
\end{itemize}

\subsection*{Lower Quarter Screen}

\begin{itemize}
  \item \textbf{Dermatomes:} L1–S2
  \item \textbf{Myotomes:} L2–S2
  \item \textbf{Reflexes:} L4 (patellar), S1 (Achilles)
  \item \textbf{MMT Examples:} DF (L4), Toe Ext (L5), PF (S1)
\end{itemize}

\begin{labactivitybox}
\textbf{Activity: Upper Quarter Screen}

\begin{itemize}
  \item Conduct dermatome light touch and compare bilaterally.
  \item Apply resistance for myotomal testing.
  \item Elicit reflexes with hammer.
  \item Record asymmetries, weakness, or absent reflexes.
\end{itemize}
\end{labactivitybox}

\begin{labactivitybox}
\textbf{Activity: Lower Quarter Screen}

\begin{itemize}
  \item Perform sensory testing for L1–S2 dermatomes.
  \item Test myotomes with standardized MMT resistance.
  \item Assess patellar and Achilles reflexes.
\end{itemize}
\end{labactivitybox}





\section*{Clinical Integration}

\begin{itemize}
\item How does cervical MMT differ in positioning, stabilization, and interpretation compared to limb testing?
\item What are the key patterns that emerge during upper and lower quarter screens when myotome, dermatome, and reflex findings are considered together?
\item How would you proceed clinically if a patient demonstrates weakness in a myotomal test but has intact sensation and reflexes?
\item In what ways do these quick screens support triage, referral, and prioritization of more detailed regional assessments?
\end{itemize}

\section*{Next Steps}

\begin{itemize}
\item Review cervical MMT positions, with attention to the smaller muscles and their segmental innervation.
\item Practice upper and lower quarter screens in a time-efficient format, emphasizing posture, patient instruction, and consistency.
\item Begin to connect MMT, reflexes, and sensation findings to root level logic and clinical decision making.
\item Preview Week 10: Skills Integration to prepare for synthesizing movement screening, neuromuscular testing, and ultrasound into a full body screen.
\end{itemize}

\section*{Checklist Summary}
\begin{itemize}
  \item [$\square$] Performed MMT for cervical flexion, extension, rotation, and sidebending.
  \item [$\square$] Identified innervation and root levels for key cervical musculature.
  \item [$\square$] Conducted the Upper Quarter Screen: myotomes, dermatomes, and reflexes.
  \item [$\square$] Conducted the Lower Quarter Screen: myotomes, dermatomes, and reflexes.
  \item [$\square$] Interpreted screen findings in the context of neuroanatomy and regional interdependence.
\end{itemize}